\documentclass{beamer}
\usetheme{Madrid}
\usepackage{amsmath,amssymb,booktabs,xcolor}
\usepackage{listings}

\lstset{
  language=Python,
  basicstyle=\ttfamily\scriptsize,
  keywordstyle=\color{blue},
  commentstyle=\color{gray},
  stringstyle=\color{orange!70!black},
  showstringspaces=false,
  breaklines=true,
  frame=single,
  backgroundcolor=\color{gray!8},
}

\title{Practical Session 5: LLM-Assisted Business Valuation}
\subtitle{Hands-On: FCFF Computation, CoT Forecasting, and Monte Carlo}
\author{Juan F.\ Imbet}
\institute{Paris Dauphine -- PSL University}
\date{Large Language Models in Finance}

\begin{document}

% ============================================================
% FRAME 1: Title
% ============================================================
\begin{frame}
\frametitle{}
\titlepage
\end{frame}

% ============================================================
% FRAME 2: Session overview
% ============================================================
\begin{frame}
\frametitle{Session Overview}

\textbf{Duration:} 1 hour

\medskip
\begin{enumerate}
  \item \textbf{Recap} (5 min) --- FCFF formula, DCF equation, error propagation, RAGAS metrics
  \item \textbf{Problem 1} (15 min) --- Compute FCFF and DCF enterprise value from raw financials
  \item \textbf{Problem 2} (10 min) --- Design a CoT forecasting prompt for a given company
  \item \textbf{Case Study} (20 min) --- Run a Monte Carlo DCF; evaluate sensitivity to WACC and $g$;
        answer three diagnostic questions
  \item \textbf{Discussion} (5 min) --- When is the DCF estimate unreliable? How would you report uncertainty?
  \item \textbf{Solutions} (5 min) --- Walkthrough of Problems 1 and 2
\end{enumerate}

\medskip
\begin{block}{What you will practise}
  Translating raw financial statement data into FCFF; structuring CoT forecasting prompts;
  parameterising and interpreting a Monte Carlo DCF; reasoning about error propagation.
\end{block}
\end{frame}

% ============================================================
% FRAME 3: Recap --- FCFF and DCF
% ============================================================
\begin{frame}
\frametitle{Recap: FCFF and DCF}

\textbf{Free Cash Flow to Firm:}
\[
  \text{FCFF}_t = \text{EBIT}_t(1-\tau) + \text{DA}_t - \text{CapEx}_t - \Delta\text{NWC}_t
\]

\textbf{DCF enterprise value:}
\[
  V = \sum_{t=1}^{n} \frac{\text{FCF}_t}{(1+\text{WACC})^t}
      + \frac{\text{FCF}_{n+1}}{(\text{WACC}-g)(1+\text{WACC})^n}
\]

\textbf{Relationship FCFF $\to$ FCFE:}
\[
  \text{FCFE}_t = \text{FCFF}_t - \text{InterestExp}_t(1-\tau) + \text{NB}_t
\]

\medskip
\textbf{Error propagation (key result):}
\[
  \frac{\partial V}{\partial g} = \frac{\text{FCF}_n(1+\text{WACC})}{(\text{WACC}-g)^2}
\]
At WACC = 9\%, $g$ = 3\%: a 100-bps error in $g$ $\to$ \textbf{15--18\% error in EV.}
\end{frame}

% ============================================================
% FRAME 4: Recap --- CoT prompting and evaluation
% ============================================================
\begin{frame}
\frametitle{Recap: CoT Forecasting and Evaluation Metrics}

\textbf{CoT four-stage structure:}
\begin{enumerate}
  \item Macro context (interest rates, GDP, inflation)
  \item Industry dynamics (drivers, headwinds, competitive dynamics)
  \item Firm-specific (competitive position, CapEx plans, margin trajectory from MD\&A)
  \item Synthesis: point estimate + P10 downside + P90 upside per year
\end{enumerate}

\medskip
\textbf{Three evaluation metrics:}
\begin{center}
\begin{tabular}{lll}
\toprule
Metric & Formula & Finance use case \\
\midrule
MAPE & $\frac{1}{n}\sum|y_t-\hat{y}_t|/|y_t|$ & Magnitude accuracy \\
RMSE & $\sqrt{\frac{1}{n}\sum(y_t-\hat{y}_t)^2}$ & Sensitivity to outliers \\
DA & $\frac{1}{n}\sum\mathbf{1}[\text{sign match}]$ & Investment decisions \\
\bottomrule
\end{tabular}
\end{center}

Always compare against random-walk baseline and analyst consensus.
\end{frame}

% ============================================================
% FRAME 5: Problem 1 --- FCFF computation
% ============================================================
\begin{frame}
\frametitle{Problem 1: Compute FCFF and Enterprise Value (15 min)}

A mid-cap industrial company reports the following for fiscal year 2023 (USD millions):

\begin{center}
\begin{tabular}{lr}
\toprule
Revenue & 2{,}850 \\
EBIT & 428 \\
Depreciation \& Amortisation & 195 \\
Capital Expenditure & 240 \\
NWC (FY2023) & 380 \\
NWC (FY2022) & 315 \\
Income Tax Expense & 85 \\
EBT (Earnings Before Tax) & 368 \\
Interest Expense & 60 \\
Net New Borrowing & +40 \\
\bottomrule
\end{tabular}
\end{center}

\textbf{Your tasks:}
\begin{enumerate}
  \item Compute FCFF using the effective tax rate (infer from income tax / EBT).
  \item Compute FCFE.
  \item Verify the FCFF-FCFE relationship.
  \item Using WACC = 10\%, 5-year explicit FCF growth = 6\%, terminal growth = 2.5\%,
        compute the DCF enterprise value. Assume the computed FCFF is the Year 0 base.
\end{enumerate}
\end{frame}

% ============================================================
% FRAME 6: Problem 2 --- CoT prompt design
% ============================================================
\begin{frame}
\frametitle{Problem 2: Design a CoT Forecasting Prompt (10 min)}

You are building a valuation agent for the same industrial company.
The MD\&A states:

\begin{quote}
\footnotesize
\textit{``We expect continued demand in our infrastructure segment driven by public
investment in grid modernisation. However, supply chain bottlenecks in copper and steel
may compress margins in the near term. We are targeting a 200-basis-point EBIT margin
improvement over three years through operational efficiency programs. CapEx will remain
elevated at 8--9\% of revenue as we expand two manufacturing facilities.''}
\end{quote}

\textbf{Your tasks:}
\begin{enumerate}
  \item Write the four-stage CoT reasoning (macro $\to$ industry $\to$ firm $\to$ synthesis)
        that a well-designed prompt should elicit.
  \item Identify at least one potential internal inconsistency in the MD\&A guidance
        that a consistency-checking step should flag.
  \item Specify the JSON output format for Year 1--3 FCF forecasts
        (base, upside, downside, justification).
\end{enumerate}
\end{frame}

% ============================================================
% FRAME 7: Case study overview
% ============================================================
\begin{frame}
\frametitle{Case Study: Monte Carlo DCF (20 min)}

\textbf{Setup:} You will parameterise and run a Monte Carlo DCF simulation
using the FCFF computed in Problem 1 as the base ($\text{FCFF}_0$).

\medskip
\textbf{Parameters:}
\begin{center}
\begin{tabular}{lcc}
\toprule
Parameter & Mean & Std Dev \\
\midrule
WACC & 10\% & 1.5\% \\
Explicit growth (per year) & 6\% & 3\% \\
Terminal growth & 2.5\% & 1\% \\
\bottomrule
\end{tabular}
\end{center}

\medskip
\textbf{Simulation spec:}
\begin{itemize}
  \item $M = 5{,}000$ draws; 5-year explicit horizon; \texttt{seed = 42}
  \item Clip terminal growth so $g_{\text{terminal}} < \text{WACC} - 0.5\%$
  \item Report: mean, median, std, P10, P25, P75, P90 of EV distribution
\end{itemize}

\medskip
\textbf{Three diagnostic questions on next slides.}
\end{frame}

% ============================================================
% FRAME 8: Case study --- step 1, code
% ============================================================
\begin{frame}[fragile]
\frametitle{Case Study Step 1: Run the Simulation}

\begin{lstlisting}
import numpy as np

def monte_carlo_dcf(fcff0, g_exp_mean, g_exp_std, wacc_mean, wacc_std,
                    g_term_mean, g_term_std, n=5, M=5000, seed=42):
    rng = np.random.default_rng(seed)
    wacc = rng.normal(wacc_mean, wacc_std, M)
    g_term = rng.normal(g_term_mean, g_term_std, M)
    g_exp = rng.normal(g_exp_mean, g_exp_std, (M, n))

    # Clip terminal growth
    g_term = np.minimum(g_term, wacc - 0.005)

    ev = np.zeros(M)
    for m in range(M):
        fcfs = [fcff0]
        for t in range(n):
            fcfs.append(fcfs[-1] * (1 + g_exp[m, t]))
        pv_fcf = sum(fcfs[t]/(1+wacc[m])**t for t in range(1, n+1))
        tv = fcfs[n]*(1+g_term[m])/(wacc[m]-g_term[m])
        ev[m] = pv_fcf + tv/(1+wacc[m])**n
    return ev

# Use FCFF0 from Problem 1
ev_sim = monte_carlo_dcf(fcff0=???, ...)
\end{lstlisting}

\textbf{Fill in \texttt{fcff0} using your answer from Problem 1.}
\end{frame}

% ============================================================
% FRAME 9: Case study --- three diagnostic questions
% ============================================================
\begin{frame}
\frametitle{Case Study: Three Diagnostic Questions}

After running the simulation, answer:

\medskip
\textbf{Question A:} Compare the P10 and P90 enterprise values.
What fraction of the total uncertainty in EV is attributable to terminal value uncertainty
vs.\ explicit-period FCF uncertainty?
(Hint: re-run with $\sigma_g = 0$ to isolate explicit-period uncertainty.)

\medskip
\textbf{Question B:} At what terminal growth rate $g$ does the median EV
approximately double relative to the base-case DCF computed in Problem 1?
(Hint: solve $TV \propto 1/(\text{WACC}-g)$ for the target $g$.)

\medskip
\textbf{Question C:} Suppose a competitor is acquired at 12$\times$ EV/EBITDA.
The company's EBITDA is EBIT + DA = $623M.
Does the comparables-based EV fall within the Monte Carlo P10--P90 range?
What does the gap (or lack of gap) tell you about market pricing?
\end{frame}

% ============================================================
% FRAME 10: Case study --- fill-in table
% ============================================================
\begin{frame}
\frametitle{Case Study: Fill In This Table}

\begin{center}
\begin{tabular}{lr}
\toprule
\textbf{Statistic} & \textbf{Value (USD M)} \\
\midrule
Mean EV & ? \\
Median EV & ? \\
Std Dev & ? \\
P10 (bear) & ? \\
P25 & ? \\
P75 & ? \\
P90 (bull) & ? \\
\midrule
P90 -- P10 range & ? \\
P90 / P10 ratio & ? \\
\midrule
Comparables EV (12× EBITDA) & ? \\
\bottomrule
\end{tabular}
\end{center}

\medskip
\textbf{Interpret:} Is the DCF estimate consistent with the comparables estimate?
What does a P90/P10 ratio $> 2$ suggest about this valuation?
\end{frame}

% ============================================================
% FRAME 11: Discussion
% ============================================================
\begin{frame}
\frametitle{Discussion: Communicating Valuation Uncertainty (5 min)}

Discuss in pairs:

\medskip
\textbf{When is the DCF estimate unreliable?}
\begin{itemize}
  \item When $g \approx \text{WACC}$ (terminal value denominator approaches zero)?
  \item When the company has negative or volatile FCF (e.g., early-stage biotech)?
  \item When the forecast horizon is very short (terminal value dominates)?
\end{itemize}

\medskip
\textbf{How would you report uncertainty to a client?}
\begin{itemize}
  \item Point estimate only? (Never --- why not?)
  \item Point estimate + range (P10--P90)?
  \item Full distribution histogram?
  \item Bull/base/bear narrative + Monte Carlo range?
\end{itemize}

\medskip
\textbf{Governance question:} You run this pipeline on 200 companies daily.
What HITL checkpoint(s) would you add, and at what threshold?
(Hint: consider P90/P10 ratio, FCFF extraction confidence, and number of non-recurring items flagged.)
\end{frame}

% ============================================================
% FRAME 12: Solution --- Problem 1
% ============================================================
\begin{frame}
\frametitle{Solutions: Problem 1 --- FCFF and DCF}

\textbf{Effective tax rate:} $\tau = 85 / 368 = 23.1\%$

\medskip
\textbf{FCFF:}
\begin{align*}
  \Delta\text{NWC} &= 380 - 315 = 65\text{ M} \\
  \text{FCFF} &= 428(1-0.231) + 195 - 240 - 65 \\
              &= 329.3 + 195 - 240 - 65 = \mathbf{219.3\text{ M}}
\end{align*}

\textbf{FCFE:}
\begin{align*}
  \text{NI} &= (428 - 60)(1-0.231) = 368 \times 0.769 = 283.0\text{ M} \\
  \text{FCFE} &= 283.0 + 195 - 240 - 65 + 40 = \mathbf{213.0\text{ M}}
\end{align*}

\textbf{Verify:} $\text{FCFF} - \text{InterestExp}(1-\tau) + \text{NB}
= 219.3 - 60(0.769) + 40 = 219.3 - 46.1 + 40 = 213.2$ M $\approx$ 213.0 M \checkmark

\medskip
\textbf{DCF:} Using python\_executor with \texttt{fcff0 = 219.3}, WACC = 10\%, $g_{\text{exp}}$ = 6\%,
$g_{\text{term}}$ = 2.5\%, $n = 5$: \textbf{EV $\approx$ 3{,}672 M}.
\end{frame}

% ============================================================
% FRAME 13: Solution --- Problem 2 highlights
% ============================================================
\begin{frame}
\frametitle{Solutions: Problem 2 --- CoT Prompt Key Points}

\textbf{Stage 3 (firm-specific) --- key reasoning points:}
\begin{itemize}
  \item Positive: grid modernisation tailwind supports infrastructure segment revenue
  \item Negative: copper/steel supply chain bottlenecks compress near-term margins
  \item Positive: 200-bps margin improvement target over 3 years
  \item Negative: CapEx at 8--9\% of revenue is elevated, limiting FCF conversion
\end{itemize}

\medskip
\textbf{Potential internal inconsistency to flag:}
\begin{quote}
\textit{``200-bps EBIT margin improvement''} while simultaneously \textit{``CapEx at 8--9\% of revenue
as we expand manufacturing facilities''}.
\end{quote}
Margin improvement is more plausible once the expansion CapEx phase ends.
If the improvements are targeted during the high-CapEx period, a consistency-checking
agent should flag this as a tension requiring clarification.

\medskip
\textbf{JSON schema:} \texttt{\{"year": 1, "base": <float>, "upside": <float>,
"downside": <float>, "justification": "<2 sentences>"\}}
\end{frame}

% ============================================================
% FRAME 14: Key takeaways
% ============================================================
\begin{frame}
\frametitle{Key Takeaways}

\begin{enumerate}
  \item \textbf{FCFF is not a data field in any filing.}
        It must be derived from EBIT, DA, CapEx, $\Delta$NWC, and the effective tax rate.
        Each input can carry LLM extraction error; errors propagate multiplicatively.

  \item \textbf{CoT separates narrative from number.}
        A good forecasting prompt forces the model to commit to a macro/industry/firm
        narrative \emph{before} generating numbers. Inconsistencies become detectable.

  \item \textbf{Terminal growth rate is the dominant uncertainty source.}
        A 100-bps error in $g$ produces a 15--18\% error in enterprise value.
        Always report Monte Carlo percentile ranges, not just point estimates.

  \item \textbf{Comparables and DCF are complements, not substitutes.}
        A large gap between the two signals either market mispricing or a flawed assumption
        in one of the models. Investigate rather than average.

  \item \textbf{Production pipelines need hallucination guards.}
        Validate every ticker; trace every financial figure to a specific accession number.
        A single hallucinated comparable can distort the entire multiple analysis.
\end{enumerate}
\end{frame}

\end{document}
